\chapter*{Apresentação}

A comunicação institucional é indispensável para a autarquia e desenvolve o papel de manter uma comunicação de duas vias, que não só transmite informações como também promove a interação entre as pessoas e dá suporte à governança na tomada de decisões, além de desenvolver e consolidar relacionamentos.

Essas políticas contribuem para as boas práticas de confiabilidade da governança, estabelecendo canais de comunicação efetivos que estimulam a interação e facilitam a interlocução com os seus públicos.

A \textbf{Política de Comunicação} do CRP~SP é resultado do \href{https://transparencia.cfp.org.br/crp06/planejamento/planejamento-estrategico-2024/}{planejamento estratégico de 2024} e cumpre o previsto na \href{https://transparencia.cfp.org.br/crp06/legislacao/regimento-interno-do-conselho-regional-de-psicologia-de-sao-paulo-6a-regiao-aprovado-pela-resolucao-no-5-de-22-de-marco-de-2023/}{Resolução CRP~SP nº~05/2023}, que dispõe sobre o regimento da autarquia.

Acredita-se que, baseados nesses aspectos, os processos de comunicação permitirão estabelecermos um canal direto entre trabalhadores, gestores, conselheiras, colaboradoras e categoria profissional que permita o fluxo de informação em todos os sentidos.

As ações de comunicação institucional a serem conduzidas serão reguladas e orientadas por esta \textbf{Política}, que relaciona as atividades a serem realizadas pela sede e subsedes.

Esta \textbf{Política}, elaborada pela Coordenadoria de Comunicação, deverá ser publicada amplamente.

De acordo com esse direcionamento, todos são responsáveis pela comunicação correta e eficaz, não sendo esta atribuição exclusiva da
Unidade de Comunicação.

É responsabilidade de todos que atuam no CRP~SP zelar pela boa imagem da instituição e cuidar para que os processos de comunicação institucional se realizem adequadamente, conforme os objetivos institucionais.

Portanto, é importante a sensibilização de todas as equipes, independentemente da área em que atuam.

A presente \textbf{Política} tem por finalidade ser um instrumento orientador e normativo de ações alinhadas aos planos estratégicos. Desta
maneira, a \textbf{Política de Comunicação} constitui um valioso instrumento a ser observado e praticado por todos, independentemente do
vínculo com o CRP~SP.

\assinatura{Conselho Regional de Psicologia de São Paulo\\XVII Plenária}